\renewcommand{\theposttitle}{What learning the guitar taught me about progress}
\renewcommand{\thepostdate}{2024-06-18}
\vspace*{\fill}
\begin{center}
\includegraphics[width=\textwidth, height=0.55\textheight, keepaspectratio]{C:/Users/James/Documents/GitHub/VisualIdeasBook/data/images/e4cdf7e00d13e62329d9b688265dc5a4.png}
\end{center}
\vspace{2mm}
\begin{center}
\begin{minipage}{0.92\textwidth}
\small
Most progress is invisible until it's not.

When I first learned to play the guitar—

E-chord? No problem. \\
C-chord? No problem. \\
But then I got to the G-chord, and I was like, what the hell is this?!

I thought, "how in the world do fingers make this shape?"

Day to day, I tried to play that freaking G-chord.

But no sound. All I felt was the pain of metal strings carving into my fingers. I couldn't see progress but for some reason I trusted the process.

And eventually one day out of no where, I heard \textit{strum}...\\
the sweet release of that G-chord! 🫠

That was the first time in my life I had proof that just because you can't see progress doesn't mean you're not making progress.
\end{minipage}
\end{center}
\vspace*{\fill}